% Options for packages loaded elsewhere
\PassOptionsToPackage{unicode}{hyperref}
\PassOptionsToPackage{hyphens}{url}
\documentclass[
]{article}
\usepackage{xcolor}
\usepackage[margin=1in]{geometry}
\usepackage{amsmath,amssymb}
\setcounter{secnumdepth}{5}
\usepackage{iftex}
\ifPDFTeX
  \usepackage[T1]{fontenc}
  \usepackage[utf8]{inputenc}
  \usepackage{textcomp} % provide euro and other symbols
\else % if luatex or xetex
  \usepackage{unicode-math} % this also loads fontspec
  \defaultfontfeatures{Scale=MatchLowercase}
  \defaultfontfeatures[\rmfamily]{Ligatures=TeX,Scale=1}
\fi
\usepackage{lmodern}
\ifPDFTeX\else
  % xetex/luatex font selection
\fi
% Use upquote if available, for straight quotes in verbatim environments
\IfFileExists{upquote.sty}{\usepackage{upquote}}{}
\IfFileExists{microtype.sty}{% use microtype if available
  \usepackage[]{microtype}
  \UseMicrotypeSet[protrusion]{basicmath} % disable protrusion for tt fonts
}{}
\makeatletter
\@ifundefined{KOMAClassName}{% if non-KOMA class
  \IfFileExists{parskip.sty}{%
    \usepackage{parskip}
  }{% else
    \setlength{\parindent}{0pt}
    \setlength{\parskip}{6pt plus 2pt minus 1pt}}
}{% if KOMA class
  \KOMAoptions{parskip=half}}
\makeatother
\usepackage{color}
\usepackage{fancyvrb}
\newcommand{\VerbBar}{|}
\newcommand{\VERB}{\Verb[commandchars=\\\{\}]}
\DefineVerbatimEnvironment{Highlighting}{Verbatim}{commandchars=\\\{\}}
% Add ',fontsize=\small' for more characters per line
\usepackage{framed}
\definecolor{shadecolor}{RGB}{248,248,248}
\newenvironment{Shaded}{\begin{snugshade}}{\end{snugshade}}
\newcommand{\AlertTok}[1]{\textcolor[rgb]{0.94,0.16,0.16}{#1}}
\newcommand{\AnnotationTok}[1]{\textcolor[rgb]{0.56,0.35,0.01}{\textbf{\textit{#1}}}}
\newcommand{\AttributeTok}[1]{\textcolor[rgb]{0.13,0.29,0.53}{#1}}
\newcommand{\BaseNTok}[1]{\textcolor[rgb]{0.00,0.00,0.81}{#1}}
\newcommand{\BuiltInTok}[1]{#1}
\newcommand{\CharTok}[1]{\textcolor[rgb]{0.31,0.60,0.02}{#1}}
\newcommand{\CommentTok}[1]{\textcolor[rgb]{0.56,0.35,0.01}{\textit{#1}}}
\newcommand{\CommentVarTok}[1]{\textcolor[rgb]{0.56,0.35,0.01}{\textbf{\textit{#1}}}}
\newcommand{\ConstantTok}[1]{\textcolor[rgb]{0.56,0.35,0.01}{#1}}
\newcommand{\ControlFlowTok}[1]{\textcolor[rgb]{0.13,0.29,0.53}{\textbf{#1}}}
\newcommand{\DataTypeTok}[1]{\textcolor[rgb]{0.13,0.29,0.53}{#1}}
\newcommand{\DecValTok}[1]{\textcolor[rgb]{0.00,0.00,0.81}{#1}}
\newcommand{\DocumentationTok}[1]{\textcolor[rgb]{0.56,0.35,0.01}{\textbf{\textit{#1}}}}
\newcommand{\ErrorTok}[1]{\textcolor[rgb]{0.64,0.00,0.00}{\textbf{#1}}}
\newcommand{\ExtensionTok}[1]{#1}
\newcommand{\FloatTok}[1]{\textcolor[rgb]{0.00,0.00,0.81}{#1}}
\newcommand{\FunctionTok}[1]{\textcolor[rgb]{0.13,0.29,0.53}{\textbf{#1}}}
\newcommand{\ImportTok}[1]{#1}
\newcommand{\InformationTok}[1]{\textcolor[rgb]{0.56,0.35,0.01}{\textbf{\textit{#1}}}}
\newcommand{\KeywordTok}[1]{\textcolor[rgb]{0.13,0.29,0.53}{\textbf{#1}}}
\newcommand{\NormalTok}[1]{#1}
\newcommand{\OperatorTok}[1]{\textcolor[rgb]{0.81,0.36,0.00}{\textbf{#1}}}
\newcommand{\OtherTok}[1]{\textcolor[rgb]{0.56,0.35,0.01}{#1}}
\newcommand{\PreprocessorTok}[1]{\textcolor[rgb]{0.56,0.35,0.01}{\textit{#1}}}
\newcommand{\RegionMarkerTok}[1]{#1}
\newcommand{\SpecialCharTok}[1]{\textcolor[rgb]{0.81,0.36,0.00}{\textbf{#1}}}
\newcommand{\SpecialStringTok}[1]{\textcolor[rgb]{0.31,0.60,0.02}{#1}}
\newcommand{\StringTok}[1]{\textcolor[rgb]{0.31,0.60,0.02}{#1}}
\newcommand{\VariableTok}[1]{\textcolor[rgb]{0.00,0.00,0.00}{#1}}
\newcommand{\VerbatimStringTok}[1]{\textcolor[rgb]{0.31,0.60,0.02}{#1}}
\newcommand{\WarningTok}[1]{\textcolor[rgb]{0.56,0.35,0.01}{\textbf{\textit{#1}}}}
\usepackage{graphicx}
\makeatletter
\newsavebox\pandoc@box
\newcommand*\pandocbounded[1]{% scales image to fit in text height/width
  \sbox\pandoc@box{#1}%
  \Gscale@div\@tempa{\textheight}{\dimexpr\ht\pandoc@box+\dp\pandoc@box\relax}%
  \Gscale@div\@tempb{\linewidth}{\wd\pandoc@box}%
  \ifdim\@tempb\p@<\@tempa\p@\let\@tempa\@tempb\fi% select the smaller of both
  \ifdim\@tempa\p@<\p@\scalebox{\@tempa}{\usebox\pandoc@box}%
  \else\usebox{\pandoc@box}%
  \fi%
}
% Set default figure placement to htbp
\def\fps@figure{htbp}
\makeatother
\setlength{\emergencystretch}{3em} % prevent overfull lines
\providecommand{\tightlist}{%
  \setlength{\itemsep}{0pt}\setlength{\parskip}{0pt}}
\usepackage{booktabs}
\usepackage{longtable}
\usepackage{array}
\usepackage{multirow}
\usepackage{wrapfig}
\usepackage{float}
\usepackage{colortbl}
\usepackage{pdflscape}
\usepackage{tabu}
\usepackage{threeparttable}
\usepackage{threeparttablex}
\usepackage[normalem]{ulem}
\usepackage{makecell}
\usepackage{xcolor}
\usepackage{bookmark}
\IfFileExists{xurl.sty}{\usepackage{xurl}}{} % add URL line breaks if available
\urlstyle{same}
\hypersetup{
  pdftitle={Machine Learning II Project Report},
  hidelinks,
  pdfcreator={LaTeX via pandoc}}

\title{Machine Learning II Project Report}
\usepackage{etoolbox}
\makeatletter
\providecommand{\subtitle}[1]{% add subtitle to \maketitle
  \apptocmd{\@title}{\par {\large #1 \par}}{}{}
}
\makeatother
\subtitle{The Analysis of the Coffee Quality Dataset}
\author{}
\date{\vspace{-2.5em}2025-12-31}

\begin{document}
\maketitle

{
\setcounter{tocdepth}{2}
\tableofcontents
}
\newpage

\section{Introduction}\label{introduction}

This report explores the Coffee Quality dataset, walking through
everything from initial descriptive analysis to the math and application
of machine learning. Our goal is to use these models to move beyond
simple numbers and extract real, meaningful insights into what defines a
high-quality cup of coffee.

\subsection{Descriptive Analysis}\label{descriptive-analysis}

\textbf{Summary of the Descriptive Analysis:} The dataset has 1,311
coffee evaluations (observations) with 44 different variables. It's a
mix of sensory scores (like how the coffee tastes) and environmental
data (like where it was grown). Out of these, \textbf{24 variables are
categorical and 20 are numerical}. For our project, we decided to use
Cupper Points as our target variable. We avoided ``Total Cup Points''
because that number is just a simple sum of the other scores, which
would make the prediction too easy and less meaningful.

\textbf{Dealing with Missing Data} One of the first things we noticed
was that a lot of the categorical (text) data is missing. Because there
were so many gaps, it was hard to use many of these variables in our
models. Instead, we focused our energy on the variables that really
mattered and tried to fill in the blanks for critical info like
Altitude.

\begin{Shaded}
\begin{Highlighting}[]
\CommentTok{\# Read the data}
\NormalTok{coffee\_data }\OtherTok{\textless{}{-}} \FunctionTok{read.csv}\NormalTok{(}\StringTok{"\textasciitilde{}/ML2 arabica{-}coffee{-}analayze/arabica\_data\_cleaned.csv"}\NormalTok{, }\AttributeTok{na.strings =} \FunctionTok{c}\NormalTok{(}\StringTok{""}\NormalTok{, }\StringTok{"NA"}\NormalTok{, }\StringTok{" "}\NormalTok{, }\StringTok{"n/a"}\NormalTok{))}
\CommentTok{\# 1. Select the features you want to focus on}
\NormalTok{selected\_features }\OtherTok{\textless{}{-}}\NormalTok{ coffee\_data }\SpecialCharTok{\%\textgreater{}\%} 
  \FunctionTok{select}\NormalTok{(Cupper.Points, altitude\_mean\_meters, Aroma, Flavor, unit\_of\_measurement, Processing.Method, }
\NormalTok{         Bag.Weight, Variety, Lot.Number, Quakers, Country.of.Origin)}

\CommentTok{\# 2. Calculate missing values and prepare for plotting}
\NormalTok{missing\_data }\OtherTok{\textless{}{-}}\NormalTok{ selected\_features }\SpecialCharTok{\%\textgreater{}\%}
  \FunctionTok{summarise}\NormalTok{(}\FunctionTok{across}\NormalTok{(}\FunctionTok{everything}\NormalTok{(), }\SpecialCharTok{\textasciitilde{}} \FunctionTok{sum}\NormalTok{(}\FunctionTok{is.na}\NormalTok{(.)))) }\SpecialCharTok{\%\textgreater{}\%}
  \FunctionTok{pivot\_longer}\NormalTok{(}\FunctionTok{everything}\NormalTok{(), }\AttributeTok{names\_to =} \StringTok{"Feature"}\NormalTok{, }\AttributeTok{values\_to =} \StringTok{"Missing\_Count"}\NormalTok{) }\SpecialCharTok{\%\textgreater{}\%}
  \FunctionTok{mutate}\NormalTok{(}\AttributeTok{Percentage =}\NormalTok{ (Missing\_Count }\SpecialCharTok{/} \FunctionTok{nrow}\NormalTok{(coffee\_data)) }\SpecialCharTok{*} \DecValTok{100}\NormalTok{) }\SpecialCharTok{\%\textgreater{}\%}
  \FunctionTok{arrange}\NormalTok{(}\FunctionTok{desc}\NormalTok{(Missing\_Count)) }\CommentTok{\# Sort so the biggest problem is first}

\CommentTok{\# 3. Create the bar plot}
\FunctionTok{ggplot}\NormalTok{(missing\_data, }\FunctionTok{aes}\NormalTok{(}\AttributeTok{x =} \FunctionTok{reorder}\NormalTok{(Feature, }\SpecialCharTok{{-}}\NormalTok{Missing\_Count), }\AttributeTok{y =}\NormalTok{ Missing\_Count)) }\SpecialCharTok{+}
  \FunctionTok{geom\_col}\NormalTok{(}\AttributeTok{fill =} \StringTok{"firebrick"}\NormalTok{, }\AttributeTok{color =} \StringTok{"white"}\NormalTok{) }\SpecialCharTok{+}
  \FunctionTok{geom\_text}\NormalTok{(}\FunctionTok{aes}\NormalTok{(}\AttributeTok{label =} \FunctionTok{paste0}\NormalTok{(}\FunctionTok{round}\NormalTok{(Percentage, }\DecValTok{1}\NormalTok{), }\StringTok{"\%"}\NormalTok{)), }
            \AttributeTok{vjust =} \SpecialCharTok{{-}}\FloatTok{0.5}\NormalTok{, }\AttributeTok{size =} \FloatTok{3.5}\NormalTok{) }\SpecialCharTok{+} \CommentTok{\# Add percentage labels above bars}
  \FunctionTok{theme\_minimal}\NormalTok{() }\SpecialCharTok{+}
  \FunctionTok{ylim}\NormalTok{(}\DecValTok{0}\NormalTok{, }\FunctionTok{max}\NormalTok{(missing\_data}\SpecialCharTok{$}\NormalTok{Missing\_Count) }\SpecialCharTok{+} \DecValTok{50}\NormalTok{) }\SpecialCharTok{+} \CommentTok{\# Add space for labels}
  \FunctionTok{labs}\NormalTok{(}\AttributeTok{title =} \StringTok{"Missing Values Count in Selected Features"}\NormalTok{,}
       \AttributeTok{x =} \StringTok{"Feature"}\NormalTok{,}
       \AttributeTok{y =} \StringTok{"Number of Missing Observations"}\NormalTok{) }\SpecialCharTok{+}
  \FunctionTok{theme}\NormalTok{(}\AttributeTok{axis.text.x =} \FunctionTok{element\_text}\NormalTok{(}\AttributeTok{angle =} \DecValTok{45}\NormalTok{, }\AttributeTok{hjust =} \DecValTok{1}\NormalTok{))}
\end{Highlighting}
\end{Shaded}

\begin{center}\includegraphics{GroupN_Report_files/figure-latex/unnamed-chunk-1-1} \end{center}

\textbf{The Distribution of the Variables} When we looked at the
histograms and boxplots for the sensory variables, we saw that they all
look very similar. Most of the scores for things like Aroma or Flavor
are clustered closely together, usually between 7 and 8. We also noticed
several outliers\emph{(coffee samples that scored much higher or lower
than the rest)}.

\begin{Shaded}
\begin{Highlighting}[]
\CommentTok{\# Prepare the data in \textquotesingle{}long\textquotesingle{} format}
\NormalTok{long\_data }\OtherTok{\textless{}{-}}\NormalTok{ coffee\_data }\SpecialCharTok{\%\textgreater{}\%}
  \FunctionTok{select}\NormalTok{(Aroma, Flavor, Cupper.Points) }\SpecialCharTok{\%\textgreater{}\%}
  \FunctionTok{pivot\_longer}\NormalTok{(}\FunctionTok{everything}\NormalTok{(), }\AttributeTok{names\_to =} \StringTok{"Feature"}\NormalTok{, }\AttributeTok{values\_to =} \StringTok{"Score"}\NormalTok{)}

\CommentTok{\# Create the Horizontal Combined Boxplot (Left side)}
\NormalTok{p\_box }\OtherTok{\textless{}{-}} \FunctionTok{ggplot}\NormalTok{(long\_data, }\FunctionTok{aes}\NormalTok{(}\AttributeTok{y =}\NormalTok{ Feature, }\AttributeTok{x =}\NormalTok{ Score, }\AttributeTok{fill =}\NormalTok{ Feature)) }\SpecialCharTok{+}
  \FunctionTok{geom\_boxplot}\NormalTok{(}\AttributeTok{alpha =} \FloatTok{0.7}\NormalTok{, }\AttributeTok{outlier.color =} \StringTok{"red"}\NormalTok{, }\AttributeTok{outlier.shape =} \DecValTok{16}\NormalTok{) }\SpecialCharTok{+}
  \FunctionTok{theme\_minimal}\NormalTok{() }\SpecialCharTok{+}
  \FunctionTok{scale\_fill\_brewer}\NormalTok{(}\AttributeTok{palette =} \StringTok{"Set2"}\NormalTok{) }\SpecialCharTok{+}
  \FunctionTok{labs}\NormalTok{(}\AttributeTok{title =} \StringTok{"Outlier \& Range Comparison"}\NormalTok{,}
       \AttributeTok{y =} \StringTok{"Feature"}\NormalTok{, }\AttributeTok{x =} \StringTok{"Score"}\NormalTok{) }\SpecialCharTok{+}
  \FunctionTok{theme}\NormalTok{(}\AttributeTok{legend.position =} \StringTok{"none"}\NormalTok{,}
        \AttributeTok{plot.title =} \FunctionTok{element\_text}\NormalTok{(}\AttributeTok{face =} \StringTok{"bold"}\NormalTok{))}

\CommentTok{\# Create Faceted Histograms (Right side)}
\NormalTok{p\_hist }\OtherTok{\textless{}{-}} \FunctionTok{ggplot}\NormalTok{(long\_data, }\FunctionTok{aes}\NormalTok{(}\AttributeTok{x =}\NormalTok{ Score, }\AttributeTok{fill =}\NormalTok{ Feature)) }\SpecialCharTok{+}
  \FunctionTok{geom\_histogram}\NormalTok{(}\AttributeTok{bins =} \DecValTok{20}\NormalTok{, }\AttributeTok{color =} \StringTok{"white"}\NormalTok{, }\AttributeTok{alpha =} \FloatTok{0.8}\NormalTok{) }\SpecialCharTok{+}
  \FunctionTok{facet\_wrap}\NormalTok{(}\SpecialCharTok{\textasciitilde{}}\NormalTok{Feature, }\AttributeTok{ncol =} \DecValTok{1}\NormalTok{) }\SpecialCharTok{+} 
  \FunctionTok{theme\_light}\NormalTok{() }\SpecialCharTok{+}
  \FunctionTok{scale\_fill\_brewer}\NormalTok{(}\AttributeTok{palette =} \StringTok{"Set2"}\NormalTok{) }\SpecialCharTok{+}
  \FunctionTok{labs}\NormalTok{(}\AttributeTok{title =} \StringTok{"Distribution Shapes"}\NormalTok{,}
       \AttributeTok{x =} \StringTok{"Score"}\NormalTok{, }\AttributeTok{y =} \StringTok{"Count"}\NormalTok{) }\SpecialCharTok{+}
  \FunctionTok{theme}\NormalTok{(}\AttributeTok{legend.position =} \StringTok{"none"}\NormalTok{,}
        \AttributeTok{strip.text =} \FunctionTok{element\_text}\NormalTok{(}\AttributeTok{size =} \DecValTok{8}\NormalTok{, }\AttributeTok{face =} \StringTok{"bold"}\NormalTok{, }\AttributeTok{color =} \StringTok{"black"}\NormalTok{),}
        \AttributeTok{plot.title =} \FunctionTok{element\_text}\NormalTok{(}\AttributeTok{face =} \StringTok{"bold"}\NormalTok{))}

\CommentTok{\# Combine Side{-}by{-}Side}
\CommentTok{\# We use \textquotesingle{}|\textquotesingle{} to put them together. }
\NormalTok{(p\_box }\SpecialCharTok{|}\NormalTok{ p\_hist) }\SpecialCharTok{+} 
  \FunctionTok{plot\_layout}\NormalTok{(}\AttributeTok{widths =} \FunctionTok{c}\NormalTok{(}\FloatTok{1.2}\NormalTok{, }\DecValTok{1}\NormalTok{)) }\SpecialCharTok{+} 
  \FunctionTok{plot\_annotation}\NormalTok{(}\AttributeTok{title =} \StringTok{"Sensory Feature Profile: Range and Distribution"}\NormalTok{,}
                  \AttributeTok{theme =} \FunctionTok{theme}\NormalTok{(}\AttributeTok{plot.title =} \FunctionTok{element\_text}\NormalTok{(}\AttributeTok{hjust =} \FloatTok{0.5}\NormalTok{, }\AttributeTok{size =} \DecValTok{16}\NormalTok{, }\AttributeTok{face =} \StringTok{"bold"}\NormalTok{)))}
\end{Highlighting}
\end{Shaded}

\begin{center}\includegraphics{GroupN_Report_files/figure-latex/unnamed-chunk-2-1} \end{center}

\textbf{Correlation Between Numerical Features} Finally, our correlation
matrix showed us how these variables lean on each other. Most sensory
features are strongly linked. if a coffee smells great (Aroma), it
almost always tastes great (Flavor). However, some things, like
Moisture, didn't follow the crowd. Moisture has a very low correlation
with the other scores and even shows a slight negative trend, meaning
more moisture sometimes leads to a slightly lower quality score.

\begin{center}\includegraphics{GroupN_Report_files/figure-latex/unnamed-chunk-3-1} \end{center}

\section{Data Preprocessing and
Cleaning}\label{data-preprocessing-and-cleaning}

\textbf{Feature selection:} In SVM Regression, categorical variables are
One-Hot Encoded using dummy variables. This means, features with many
unique values will lead to a sparse feature matrix. This is
disadvantageous as it increases computational cost and poses the risk of
overfitting. Additionally, categories with only a few or even only one
observation most likely will not add to model robustness, but are adding
noise to the signal. Therefore, we are removing categorical features
with many unique values, only keeping columns with a low number of
unique values. E.g. we are not including the ``Region'' column with 344
unique values in the model training, because we'd on average have less
than 4 observations per Region. For features with a high number of
missing values, we either remove the column or impute the missing
values. For instance for the ``Variety'' column, we have about 15\%
missing values and 8\% entries ``other'' at 30 different varieties.
Keeping those, would mean adding 30 sparse columns, and at the same time
the information gain is most likely not too high with this many missing
values. Therefore, we are not including this feature.

\textbf{Imputation:} For the altitude column, we have about 17\% missing
values. Here, we are imputing the missing values from the ``Region''
column. For each observation with missing altitude, we are imputing the
median altitude of all coffees from the same region. This covers almost
all missing values. For the remaining observations, we are imputing the
mean altitude of all coffees from the same country of origin. As the
``Region'' and ``Country'' columns are not used in the model training,
we are not creating correlated features with this imputation. We are
imputing the missing values after the train/test split, using only
values from the training set, to avoid data leakage from the test set
into the training set.

\textbf{Further preprocessing:} The bag weight column contains both
pounds and kilograms, so we are converting all to kilograms. However,
for some entries, it is unclear which unit is used. For these, we are
averaging the weight in pounds and kilograms. This means for those
entries the weight is certainly incorrect, but should be within an order
of magnitude of the correct weight. If we can assume that speciality
coffee is sold in smaller quantities, we can still gain valuable
information from the weight.

\textbf{Scaling:} The SVM method used for model training has a built in
scaling functionality. Both predictor and target variables are scaled to
have expected value 0 and variance 1.
{[}\url{https://rdrr.io/rforge/e1071/man/svm.html}{]} This
transformation is based on the training data and applied to train and
test data sets.

\section{A Brief Mathematical Overview of the ML Methods
Used}\label{a-brief-mathematical-overview-of-the-ml-methods-used}

\subsection{Support Vector Machine for
Regression}\label{support-vector-machine-for-regression}

\begin{Shaded}
\begin{Highlighting}[]
\FunctionTok{set.seed}\NormalTok{(}\DecValTok{123}\NormalTok{)}


\NormalTok{x }\OtherTok{\textless{}{-}} \FunctionTok{seq}\NormalTok{(}\SpecialCharTok{{-}}\DecValTok{10}\NormalTok{, }\DecValTok{10}\NormalTok{, }\FloatTok{0.5}\NormalTok{) }\SpecialCharTok{+} \FunctionTok{rnorm}\NormalTok{(}\FunctionTok{length}\NormalTok{(}\FunctionTok{seq}\NormalTok{(}\SpecialCharTok{{-}}\DecValTok{10}\NormalTok{, }\DecValTok{10}\NormalTok{, }\FloatTok{0.5}\NormalTok{)), }\AttributeTok{mean =} \DecValTok{0}\NormalTok{, }\AttributeTok{sd =} \FloatTok{0.1}\NormalTok{)}
\NormalTok{y }\OtherTok{\textless{}{-}} \DecValTok{1}\SpecialCharTok{/}\DecValTok{170} \SpecialCharTok{*}\NormalTok{ (x}\SpecialCharTok{\^{}}\DecValTok{3} \SpecialCharTok{{-}} \DecValTok{3}\SpecialCharTok{*}\NormalTok{ x}\SpecialCharTok{\^{}}\DecValTok{2} \SpecialCharTok{{-}} \DecValTok{70}\SpecialCharTok{*}\NormalTok{x }\SpecialCharTok{+} \DecValTok{5} \SpecialCharTok{+} \FunctionTok{rnorm}\NormalTok{(}\FunctionTok{length}\NormalTok{(x), }\AttributeTok{mean =} \DecValTok{0}\NormalTok{, }\AttributeTok{sd =} \DecValTok{170}\NormalTok{))}

\FunctionTok{library}\NormalTok{(e1071)}

\NormalTok{epsilon }\OtherTok{\textless{}{-}} \DecValTok{1}
\NormalTok{cost }\OtherTok{\textless{}{-}} \DecValTok{1}

\NormalTok{model }\OtherTok{\textless{}{-}} \FunctionTok{svm}\NormalTok{(y }\SpecialCharTok{\textasciitilde{}}\NormalTok{ x, }\AttributeTok{data =} \FunctionTok{data.frame}\NormalTok{(}\AttributeTok{x =}\NormalTok{ x, }\AttributeTok{y =}\NormalTok{ y), }\AttributeTok{epsilon =}\NormalTok{ epsilon, }\AttributeTok{cost =}\NormalTok{ cost)}
\NormalTok{predictedY }\OtherTok{\textless{}{-}} \FunctionTok{predict}\NormalTok{(model, }\FunctionTok{data.frame}\NormalTok{(}\AttributeTok{x =}\NormalTok{ x))}

\FunctionTok{plot}\NormalTok{(x, y, }\AttributeTok{pch =} \DecValTok{16}\NormalTok{, }\AttributeTok{col =} \StringTok{"grey70"}\NormalTok{)}
\FunctionTok{lines}\NormalTok{(x, predictedY, }\AttributeTok{col =} \StringTok{"darkgreen"}\NormalTok{)}


\CommentTok{\# Add ε{-}tube}
\FunctionTok{polygon}\NormalTok{(}
 \FunctionTok{c}\NormalTok{(x, }\FunctionTok{rev}\NormalTok{(x)),}
 \FunctionTok{c}\NormalTok{(predictedY }\SpecialCharTok{+}\NormalTok{ epsilon, }\FunctionTok{rev}\NormalTok{(predictedY }\SpecialCharTok{{-}}\NormalTok{ epsilon)),}
 \AttributeTok{border =} \ConstantTok{NA}\NormalTok{,}
 \AttributeTok{col =} \FunctionTok{rgb}\NormalTok{(}\DecValTok{1}\NormalTok{, }\DecValTok{0}\NormalTok{, }\DecValTok{0}\NormalTok{, }\FloatTok{0.15}\NormalTok{)}
\NormalTok{)}

\CommentTok{\# mark the 17 points that are the furthest from the predictedY }
\NormalTok{idx\_furthest }\OtherTok{\textless{}{-}} \FunctionTok{order}\NormalTok{(}\FunctionTok{abs}\NormalTok{(predictedY }\SpecialCharTok{{-}}\NormalTok{ y), }\AttributeTok{decreasing =} \ConstantTok{TRUE}\NormalTok{)[}\DecValTok{1}\SpecialCharTok{:}\DecValTok{17}\NormalTok{]}
\FunctionTok{points}\NormalTok{(x[idx\_furthest], y[idx\_furthest], }
       \AttributeTok{col =} \StringTok{"\#000000"}\NormalTok{, }
       \AttributeTok{bg =} \StringTok{"\#ffffff00"}\NormalTok{,}
       \AttributeTok{pch =} \DecValTok{21}\NormalTok{,}
       \AttributeTok{cex =} \FloatTok{1.6}\NormalTok{,}
       \AttributeTok{lwd =} \DecValTok{1}\NormalTok{)}


\CommentTok{\# Add an arrow (dashed, double sided, between ypred at x = 10 and ypred at x = 10   + 1 )}
\FunctionTok{arrows}\NormalTok{(}\DecValTok{10}\NormalTok{, predictedY[}\FunctionTok{length}\NormalTok{(predictedY)], }\DecValTok{10}\NormalTok{, predictedY[}\FunctionTok{length}\NormalTok{(predictedY)] }\SpecialCharTok{+}\NormalTok{ epsilon, }\AttributeTok{length =} \FloatTok{0.1}\NormalTok{, }\AttributeTok{lty =} \DecValTok{2}\NormalTok{, }\AttributeTok{col =} \StringTok{"black"}\NormalTok{, }\AttributeTok{code =} \DecValTok{3}\NormalTok{)}
\CommentTok{\# add a text "ε" at the end of the arrow}
\FunctionTok{text}\NormalTok{(}\DecValTok{10}\NormalTok{, predictedY[}\FunctionTok{length}\NormalTok{(predictedY)] }\SpecialCharTok{+} \DecValTok{1}\SpecialCharTok{/}\DecValTok{2}\NormalTok{, }\FunctionTok{expression}\NormalTok{(epsilon), }\AttributeTok{pos =} \DecValTok{4}\NormalTok{)}

\CommentTok{\# Add an arrow between 5,{-}5 and 5,{-}2.5 with text xsi\_i star (flipped direction)}
\FunctionTok{arrows}\NormalTok{(}\FloatTok{5.05}\NormalTok{, }\SpecialCharTok{{-}}\DecValTok{4}\NormalTok{, }\FloatTok{5.05}\NormalTok{, }\SpecialCharTok{{-}}\FloatTok{2.28}\NormalTok{, }\AttributeTok{length =} \FloatTok{0.1}\NormalTok{, }\AttributeTok{lty =} \DecValTok{2}\NormalTok{, }\AttributeTok{col =} \StringTok{"black"}\NormalTok{)}
\FunctionTok{text}\NormalTok{(}\FloatTok{5.05}\NormalTok{, }\SpecialCharTok{{-}}\FloatTok{3.2}\NormalTok{, }\FunctionTok{expression}\NormalTok{(xi[j]}\SpecialCharTok{\^{}}\StringTok{"*"}\NormalTok{), }\AttributeTok{pos =} \DecValTok{4}\NormalTok{)}


\CommentTok{\# Add an arrow between the 13th last point and the 13th last point {-} 1 with text xsi\_j}
\FunctionTok{arrows}\NormalTok{(x[}\FunctionTok{length}\NormalTok{(x) }\SpecialCharTok{{-}} \DecValTok{12}\NormalTok{], y[}\FunctionTok{length}\NormalTok{(y) }\SpecialCharTok{{-}} \DecValTok{12}\NormalTok{], x[}\FunctionTok{length}\NormalTok{(x) }\SpecialCharTok{{-}} \DecValTok{12}\NormalTok{], y[}\FunctionTok{length}\NormalTok{(y) }\SpecialCharTok{{-}} \DecValTok{12}\NormalTok{]}\SpecialCharTok{{-}} \FloatTok{0.5}\NormalTok{, }\AttributeTok{length =} \FloatTok{0.1}\NormalTok{, }\AttributeTok{lty =} \DecValTok{2}\NormalTok{, }\AttributeTok{col =} \StringTok{"black"}\NormalTok{)}
\FunctionTok{text}\NormalTok{(x[}\FunctionTok{length}\NormalTok{(x) }\SpecialCharTok{{-}} \DecValTok{12}\NormalTok{], y[}\FunctionTok{length}\NormalTok{(y) }\SpecialCharTok{{-}} \DecValTok{12}\NormalTok{] }\SpecialCharTok{{-}} \FloatTok{0.3}\NormalTok{, }\FunctionTok{expression}\NormalTok{(xi[k]), }\AttributeTok{pos =} \DecValTok{4}\NormalTok{)}
\end{Highlighting}
\end{Shaded}

\begin{center}\includegraphics[width=1\linewidth]{GroupN_Report_files/figure-latex/unnamed-chunk-4-1} \end{center}

Support Vector Regression expands support Vector Machines from discrete
classification to a continuos regression model.

For the regression task, instead of fitting a hyperplane, we are trying
to fit a (multidimensional) function \(f(x)\) to the data. Similar as in
the classification task, we draw a margin around the function. However,
now the margin \(\varepsilon\) is fixed at the start. (Note that
\(\varepsilon\) here, is not the slack variable as in Hasties and
Tibshirani, but rather defines the width of the margin. Also, we don't
aim to optimize the margin anymore, but instead fix it before the
computation starts). We want to find a function f(x) that is as flat as
possible, while fitting all points into the required margin. Again, we
allow for some slack using slack variables. A wider margin decreases the
risk of overfitting, while a smaller margin captures more intricacies of
the data.

Mathematically speaking, want to find a function \(f(x)\) with
\(| f(x_i) - y_i |  \leq \varepsilon\) for all \(i\). To avoid
overfitting, we also introduce the constraint of making the function as
flat as possible. ``Flatness'' in this context is a measure of how
sensitive the function is to change. For a linear function
\(f(x) = x^\top \beta + b\) increasing this flatness can be expressed
through minimizing the norm of the function gradient:
\(\lVert \beta \rVert\). Minimizing the expression
\(\frac{1}{2} \lVert \beta \rVert^2\) leads to the same optimum, while
allowing for more elegant mathematical solutions, and is therefore used
for computation.

Again, we want to allow for some slack, so as in the SVM for
classification, we are introducing slack variables \(\xi_i\) and
\(\xi_i^\ast\). These slack variables allow for some room for margin
violations. The optimization problem we want to solve is to minimize:

\[
\frac{1}{2} \lVert \beta \rVert^2 + K \sum_{i=1}^n (\xi_i + \xi_i^\ast)
\] subject to: \[
y_i - f(x_i) \leq \varepsilon + \xi_i, 
\] \[
f(x_i) - y_i \leq \varepsilon + \xi_i^\ast
\] and \[
\xi_i > 0, \xi_i^\ast > 0 \quad \forall i
\]

where \(K\) is fixed and \(\beta\), \(\xi\) and \(\xi^\ast\) are
variable. The factor \(K\) is introduced as cost for regularization
purposes. It is choosen a priori to determine how strong the tradeoff
between flatness and overfitting should be. A larger \(K\) leads to a
stronger impact of the term, therefore penalizing stronger margin
violations more heavily and resulting in a less flat function curve
(higher bias, lower variance). Conversly, a smaller \(K\) leads to a
flatter curve (lower bias, higher variance).

As in the SVM for classification, this can be generalized to non-linear
functions by transforming the features using kernels.

\subsection{XGBoost}\label{xgboost}

\section{Application of SVM Method}\label{application-of-svm-method}

\subsubsection{Support Vector Regression - Hyperparameter
Optimization}\label{support-vector-regression---hyperparameter-optimization}

For the prediction of cupper points we train an SVM for regression on
the data set, holding out the test data set for final evaluation. To
choose appropriate hyperparameters, we conduct a hyperparameter
optimization on the training data set using 10-fold cross-validation. We
use a grid search for different values for cost and epsilon for a
radial, linear, and 2-degree polynomial kernel. We compare the mean
absolute error between runs to choose the best-performing hyperparameter
combination. Finally, we train a model with these hyperparameters on the
entire training set and evaluate on the test data set.

\begin{center}\includegraphics{GroupN_Report_files/figure-latex/unnamed-chunk-7-1} \end{center}

These figures show the mean absolute error for the different
hyperparameters by kernel. The best performing hyperparameter
combination for each kernel are marked with orange points.

\begin{Shaded}
\begin{Highlighting}[]
\CommentTok{\# Compare best results for each kernel by getting the row with the lowest error}
\NormalTok{best\_runs }\OtherTok{\textless{}{-}}\NormalTok{ best\_runs[}\FunctionTok{order}\NormalTok{(best\_runs}\SpecialCharTok{$}\NormalTok{err),]}

\NormalTok{best\_run }\OtherTok{\textless{}{-}}\NormalTok{ best\_runs[}\DecValTok{1}\NormalTok{,]}

\FunctionTok{cat}\NormalTok{(}\StringTok{"Best SVM kernel and parameters selected:}\SpecialCharTok{\textbackslash{}n}\StringTok{"}\NormalTok{)}
\end{Highlighting}
\end{Shaded}

\begin{verbatim}
## Best SVM kernel and parameters selected:
\end{verbatim}

\begin{Shaded}
\begin{Highlighting}[]
\FunctionTok{print}\NormalTok{(best\_run[,}\FunctionTok{c}\NormalTok{(}\StringTok{"kernel"}\NormalTok{, }\StringTok{"epsilon"}\NormalTok{, }\StringTok{"cost"}\NormalTok{)])}
\end{Highlighting}
\end{Shaded}

\begin{verbatim}
##   kernel epsilon      cost
## 2 linear     0.1 0.0078125
\end{verbatim}

\begin{Shaded}
\begin{Highlighting}[]
\CommentTok{\# Train final model with best parameters}
\NormalTok{epsilon }\OtherTok{=}\NormalTok{ best\_run}\SpecialCharTok{$}\NormalTok{epsilon}
\NormalTok{cost }\OtherTok{=}\NormalTok{ best\_run}\SpecialCharTok{$}\NormalTok{cost}

\NormalTok{svm\_model }\OtherTok{\textless{}{-}} \FunctionTok{svm}\NormalTok{(Cupper.Points }\SpecialCharTok{\textasciitilde{}}\NormalTok{ ., }\AttributeTok{data =}\NormalTok{ train\_data,}
                 \AttributeTok{cost =}\NormalTok{ cost,}
                 \AttributeTok{epsilon =}\NormalTok{ epsilon, }
                 \AttributeTok{kernel =}\NormalTok{ best\_run}\SpecialCharTok{$}\NormalTok{kernel)}

\CommentTok{\# Predict on test set}
\NormalTok{predictions }\OtherTok{\textless{}{-}} \FunctionTok{predict}\NormalTok{(svm\_model, test\_data)}

\CommentTok{\# Evaluate model performance}
\NormalTok{mae }\OtherTok{\textless{}{-}} \FunctionTok{mean}\NormalTok{(}\FunctionTok{abs}\NormalTok{(predictions }\SpecialCharTok{{-}}\NormalTok{ test\_data[, }\StringTok{"Cupper.Points"}\NormalTok{]), }\AttributeTok{na.rm =} \ConstantTok{TRUE}\NormalTok{)}
\NormalTok{mse }\OtherTok{\textless{}{-}} \FunctionTok{mean}\NormalTok{((predictions }\SpecialCharTok{{-}}\NormalTok{ test\_data[, }\StringTok{"Cupper.Points"}\NormalTok{])}\SpecialCharTok{\^{}}\DecValTok{2}\NormalTok{, }\AttributeTok{na.rm =} \ConstantTok{TRUE}\NormalTok{)}
\NormalTok{rmse }\OtherTok{\textless{}{-}} \FunctionTok{sqrt}\NormalTok{(mse)}
\CommentTok{\# print(paste("RMSE:", rmse))}
\FunctionTok{print}\NormalTok{(}\FunctionTok{paste}\NormalTok{(}\StringTok{"MAE:"}\NormalTok{, mae))}
\end{Highlighting}
\end{Shaded}

\begin{verbatim}
## [1] "MAE: 0.143099275160558"
\end{verbatim}

The best performing kernel is the linear kernel, which means the
relationship between the features and the target variable is best
approximated by a multidimensional hyperplane.

\begin{center}\includegraphics{GroupN_Report_files/figure-latex/unnamed-chunk-9-1} \end{center}

Here, we can see the predictions vs.~the true values for the test data
set. We can see that aside from two outliers, the predictions are close
to the true values.

\section{Application of XGBoost
Method}\label{application-of-xgboost-method}

\end{document}
